Let $E$ be the selected element width in bytes, $N=VLEN/(8E)$, and $G$ be
the typed governing-predicate image whose significant bits are
\texttt{Pg[iE]} for $0\leq i<N$.  \texttt{Pg} and \texttt{Ps} shall name
different (:term:VECTOR.terms.predicate_register|plural:).  Before attempting a lane, VGATHER replaces
\texttt{Ps} by its typed image \texttt{Ps AND G}; all bits outside the typed
predicate image are cleared.  Thus completed lanes that are no longer governed are discarded.

If \texttt{G AND NOT Ps} is empty, the instruction completes normally with
\texttt{Ps == G}.  Otherwise it selects one pending lane $i$, calculates that
lane's (:term:base.terms.effective_address:) from the selected address form, and issues exactly one $E$-byte
normal-memory load through the default data segment.  On success it replaces
only destination lane $i$, sets \texttt{Ps[iE]}, and commits those effects at
a lane-completion boundary.  If more lanes remain pending, the PC continues to
identify this VGATHER and execution may admit an asynchronous event before
another lane is attempted.  A lane-completion boundary is not an instruction
retirement or a debug-trace boundary.

A synchronous memory fault commits no destination value and does not set the
faulting lane's completion bit; the typed normalization of \texttt{Ps} remains
committed, so retrying at the saved PC does not repeat an already completed
lane.  Software emulation of a faulting lane shall write the destination lane
before setting its completion bit.  The vector address and destination
operands shall not name the same (:term:VECTOR.terms.vector_register:).  VGATHER preserves FLAGS and
is not permitted for MMIO.
